% !PRINTED 695
\weberpage{695}
%<<CONT
is obtained, and if one here sets $m=p$, it follows that there are
classes of order $p$ of the second kind.

\begin{center}
§ 199. \gesperrt{The Riemann-Roch Theorem.}
\end{center}

The theorem known under the name of the \gesperrt{Riemann-Roch} theorem
has the purpose of determining the dimensions of polygon
 classes.

We consider first a \gesperrt{proper polygon class}
 $A$ of order $n$, take in it two coprime
polygons $\mathfrak{A}$, $\mathfrak{A}'$ and set:
\[
z = \frac{\mathfrak{A}'}{\mathfrak{A}}.
\]

If then $\mathfrak{A}''$ is another polygon of the same class, we
set
\[
\omega = \frac{\mathfrak{A}''}{\mathfrak{A}},\quad \frac{\omega}{z} = \frac{\mathfrak{A}''}{\mathfrak{A}'}.
\]

Thus $\omega$ is an integral function of $z$, and $\omega:z$ an
integral function of $1:z$. If then $\omega$ is not constant, the
exponent $r$ of $\omega$ is equal to 1. Conversely, every integral
function $\omega$ of $z$ of exponent 0 or 1 is likewise representable in the form
$\mathfrak{A}'':\mathfrak{A}$.

We thus obtain the whole class $A$ from the totality
of the integral functions of $z$ whose exponent is $\leqq 1$. These
we obtain easily from the normal basis $(\lambda_1, \lambda_2, \ldots, \lambda_n)$ in $z$,
with the sequence of exponents $r_1, r_2, \ldots, r_n$ (§ 193). If $\lambda_s$ is the
last of these functions whose exponent is equal to 1 or 0, then
every integral function whose exponent is $\leqq 1$ is representable in the form

\weq{1}{\omega = c_0 z + c_1\lambda_1 + c_2\lambda_2 + \cdots + c_s\lambda_s}

($c_0$ must be constant, because by the definition
of the normal basis it must be an integral function, and if its
exponent were positive, the exponent of $\omega$ would be greater than
one). Accordingly

\weq{2}{(O, A) = s + 1,}

and this number is therefore always $\leqq n + 1$.

The upper limit $n+1$ is reached only when all
$r_2, r_3, \ldots, r_n = 1$ and consequently $p = 0$.

Since at least two polygons are contained in a proper class,
so that $(O, A) \geqq 2$ must hold, $n = 1$ can occur only
in the case $p = 0$.
